% ===== ch07 : 原文 main.tex 第 1805--2169 行（5a G-等变层）=====
% 本文件是原文 LaTeX 源码的逐行分片（去掉整段注释掉的旧稿），仅供翻译时对照；最终译文写到 src/chunks/ch07.tex。
\section{$G$-equivariant sheaves\label{sec:G-equiv-sheaves}}
In this section we recall the definition of $G$-equivariant sheaves (from e.g. \cite[\href{https://stacks.math.columbia.edu/tag/03LE}{Tag 03LE}]{stacks-project}) and establish the relation between $G$-equivariant coherent sheaves/vector bundles over a scheme $X$ and $H$-equivariant coherent sheaves/vector bundles over $X$, with $H\subseteq G$ a subgroup.

\subsection{Categories of $G$-equivariant $\mathcal{O}_X$-modules.}
Let $G=(G, m, e)$ a group scheme over $\field$ with multiplication $m : G \times G \rightarrow G$ and identity section $e : \Spec(\field) \rightarrow G$,
and
$X$ a scheme over $\field$
equipped with a $G$-action
$\sigma : G \times X \rightarrow X$. See \cite[\href{https://stacks.math.columbia.edu/tag/022R}{Section 022R}]{stacks-project} and \cite[\href{https://stacks.math.columbia.edu/tag/01II}{Section 01II}]{stacks-project} for some of the basic definitions.

\begin{defn}\label{defn:EquivariantSheaves}
A \emph{$G$-equivariant $\mathcal{O}_X$-module}
$(\mathcal{F}, \theta)$
is an $\mathcal{O}_X$-module $\mathcal{F}$
together with an isomorphism
% ---- 原文行 1820 ----
$$
\theta : \sigma^* \mathcal{F} \xrightarrow{\cong} \left(\pr_X^{G \times X}\right)^* \mathcal{F}
$$
of $\mathcal{O}_{G\times X}$-modules
(called a \emph{$G$-linearization} of $\mathcal{F}$),
where $\pr_X^{G \times X} : G \times X \rightarrow X$ is the canonical projection,
such that the following conditions hold:
\begin{itemize}
\item \emph{(cocycle condition)}
$
(m \times \id_X)^* \theta
=
(\id_G \times \sigma)^* \theta \comp \left(\pr_{G \times X}^{G \times G \times X}\right)^* \theta.
$

% Equivalently,
% the following diagram commutes in $\Mod(G\times G\times X)$:
% $$
% \begin{tikzcd}[row sep=3em, column sep={3.5em, 1em}]
% (\id_G \times \sigma)^* \sigma^* \mathcal{F} 
% ---- 原文行 1840 ----
%   \arrow[r, "{(\id_G \times \sigma)^*\theta}"] 
%   \arrow[d, equal] 
% &
% (\id_G \times \sigma)^* \left(\pr_X^{G \times X}\right)^* \mathcal{F}
%   \arrow[r, equal]
% &[-3em]
% \left(\pr_{G \times X}^{G \times G \times X}\right)^* \sigma^* \mathcal{F} 
%   \arrow[d, "{\left(\pr_{G \times X}^{G \times G \times X}\right)^*\theta}"] \\
% (m \times \id_X)^* \sigma^* \mathcal{F} 
%   \arrow[r, "{(m \times \id_X)^*\theta}"] 
% &
% (m \times \id_X)^* \left(\pr_X^{G \times X}\right)^* \mathcal{F}
%   \arrow[r, equal]
% &
% \left(\pr_{G \times X}^{G \times G \times X}\right)^* \left(\pr_X^{G \times X}\right)^* \mathcal{F}
% \end{tikzcd}
% $$
\item \emph{(unit condition)}
$(e \times \id_X)^*\theta = \id_{\mathcal{F}}$.

% ---- 原文行 1860 ----
% Equivalently,
% the following diagram commutes in $\Mod(X)$:
% $$
% \begin{tikzcd}[row sep=2.5em, column sep=4em]
% (e \times \id_X)^* \sigma^* \mathcal{F} 
%   \arrow[r, "{(e \times \id_X)^*\theta}"] 
%   \arrow[d, equal] 
% & (e \times \id_X)^* \left(\pr_X^{G \times X}\right)^* \mathcal{F} 
%   \arrow[d, equal] \\
% \mathcal{F} 
%   \arrow[r, "\id_{\mathcal{F}}"] 
% & \mathcal{F}
% \end{tikzcd}
% $$
\end{itemize}

A \emph{morphism of $G$-equivariant $\mathcal{O}_X$-modules}
from
$\mathcal{F}=(\mathcal{F}, \theta)$
to
% ---- 原文行 1880 ----
$\mathcal{F}'=(\mathcal{F}', \theta')$
is a morphism
$
f : \mathcal{F} \rightarrow \mathcal{F}'
$
of $\mathcal{O}_X$-modules
that commutes with the $G$-linearizations,
i.e.,
$$
\sigma^* f \comp \theta'
=
\theta \comp \left(\pr_X^{G \times X}\right)^* f
$$
as morphisms of $\mathcal{O}_{G \times X}$-modules.
% Equivalently,
% the following diagram commutes in $\Mod(G\times X)$:
% $$
% \begin{tikzcd}[row sep=large, column sep=large]
% \sigma^* \mathcal{F} \arrow[r, "\sigma^* f"]
%   \arrow[d, "\theta"']
% ---- 原文行 1900 ----
% & 
% \sigma^* \mathcal{F}' \arrow[d, "\theta'"]
% \\
% \left(\pr_X^{G \times X}\right)^* \mathcal{F}
%   \arrow[r, "(\pr_X^{G \times X})^* f"]
% &
% \left(\pr_X^{G \times X}\right)^* \mathcal{F}'
% \end{tikzcd}
% $$



%$$
%\Hom_{\Mod_G(X)}\big((\mathcal{F}, \theta), \; (\mathcal{F}', \theta')\big)
%:=
%\left\{ f \in \Hom_{\Mod(X)}(\mathcal{F}, \mathcal{F}') \;\middle|\; \sigma^* f \comp \theta' = \theta \comp \left(\pr_X^{G \times X}\right)^* f \right\}
%$$


We denote by
% ---- 原文行 1920 ----
$\Mod_G(X)$
the \emph{category of $G$-equivariant $\mathcal{O}_X$-modules}.
A $G$-equivariant $\mathcal{O}_X$-module $(\mathcal{F},\theta)$
is said to be
\emph{quasi-coherent}, \emph{coherent},
or \emph{locally free of finite rank}
if $\mathcal{F}$
is an $\mathcal{O}_X$-module having the corresponding property.
We denote by
$\QCoh_G(X)$,
$\Coh_G(X)$,
and
$\VB_G(X)$
the full subcategories of $\Mod_G(X)$
whose objects are
$G$-equivariant quasi-coherent, coherent, and finite rank locally free $\mathcal{O}_X$-modules,
respectively,
so that they satisfy
$$
\VB_G(X)\subset\Coh_G(X)\subset\QCoh_G(X)\subset\Mod_G(X).
% ---- 原文行 1940 ----
$$
\end{defn}




\begin{defn}[Internal tensor product of equivariant sheaves]
\label{defn:TensorProductOfEquivariantModules}
Let $(\mathcal{F}, \theta)$ and $(\mathcal{F}', \theta')$ be two $G$-equivariant $\mathcal{O}_X$-modules in $\Mod_G(X)$.
The \emph{(internal) tensor product} of $(\mathcal{F}, \theta)$ and $(\mathcal{F}', \theta')$
$$
(\mathcal{F}, \theta)
\underset{\mathcal{O}_X}{\otimes}
(\mathcal{F}', \theta')
:=
\left(\mathcal{F}\underset{\mathcal{O}_X}{\otimes}\mathcal{F}', \, \theta\otimes\theta'\right)
\in
\Mod_G(X),
$$
is defined as the $\mathcal{O}_X$-module $\mathcal{F}\underset{\mathcal{O}_X}{\otimes}\mathcal{F}'$
% ---- 原文行 1960 ----
together with the $G$-linearization
$$
\theta\otimes\theta'
:
\sigma^*\left(
\mathcal{F}
\underset{\mathcal{O}_X}{\otimes}
\mathcal{F}'
\right)
\xrightarrow{\cong}
\left(\pr_X^{G \times X}\right)^*
\left(\mathcal{F}
\underset{\mathcal{O}_X}{\otimes}
\mathcal{F}'\right)
$$
given by the tensor product $\theta \otimes \theta'$
under the canonical isomorphisms
$$
\sigma^*\left(\mathcal{F}\underset{\mathcal{O}_X}{\otimes}\mathcal{F}'\right)
\cong
% ---- 原文行 1980 ----
\sigma^* \mathcal{F} \underset{\mathcal{O}_{G \times X}}{\otimes} \sigma^* \mathcal{F}'
\quad\text{and}\quad
\left(\pr_X^{G \times X}\right)^* \left(\mathcal{F}\underset{\mathcal{O}_X}{\otimes}\mathcal{F}'\right)
\cong
\left(\pr_X^{G \times X}\right)^*\mathcal{F} \underset{\mathcal{O}_{G \times X}}{\otimes} \left(\pr_X^{G \times X}\right)^*\mathcal{F}'.
$$


The \emph{unit object} for the tensor product is defined by
$(\mathcal{O}_X, \id_{\mathcal{O}_{G\times X}}) \in \Mod_G(X)$,
where
$
\id_{\mathcal{O}_{G\times X}}
:
\sigma^* \mathcal{O}_X
\xrightarrow{\cong}
\left(\pr_X^{G \times X}\right)^* \mathcal{O}_X
$
is induced from the canonical isomorphisms of sheaves
$
% ---- 原文行 2000 ----
\sigma^* \mathcal{O}_X
\cong
\mathcal{O}_{G\times X}
$
and
$
\left(\pr_X^{G \times X}\right)^* \mathcal{O}_X
\cong
\mathcal{O}_{G\times X}.
$


Equipped with this tensor product,
$\Mod_G(X)$
forms a
\emph{symmetric monoidal category}
$
\left(
\Mod_G(X),
\underset{\mathcal{O}_X}{\otimes},
% ---- 原文行 2020 ----
\left(\mathcal{O}_X,\theta(0)\right)
\right).
$
Furthermore, the operation $\underset{\mathcal{O}_X}{\otimes}$ restricts to endow
$\QCoh_G(X)$, $\Coh_G(X)$, and $\VB_G(X)$
with symmetric monoidal category structures as well.

\end{defn}



\begin{defn}[Internal hom of equivariant sheaves]
\label{defn:InternalHomOfEquivariantModules}
Let $(\mathcal{F}, \theta)$ and $(\mathcal{F}', \theta')$ be two $G$-equivariant $\mathcal{O}_X$-modules in $\Mod_G(X)$
such that
$\mathcal{F}$ is finitely presented\footnote{
An $\mathcal{O}_X$-module $\mathcal{F}\in\Mod(X)$
is finitely presented if
(1)
$\mathcal{F}\in\VB(X)$ is finite rank locally free,
% ---- 原文行 2040 ----
or
(2)
$X$ is locally Noetherian and $\mathcal{F}\in\Coh(X)$ is a coherent sheaf.
}.
The \emph{internal hom} of $(\mathcal{F}, \theta)$ and $(\mathcal{F}', \theta')$
$$
\mathcal{H}om_{\mathcal{O}_X}\left( (\mathcal{F}, \theta), (\mathcal{F}', \theta') \right)
:=
\left( \mathcal{H}om_{\mathcal{O}_X}(\mathcal{F}, \mathcal{F}'), \, \Theta \right)
\in
\Mod_G(X),
$$
is defined to be the $\mathcal{O}_X$-module $\mathcal{H}om_{\mathcal{O}_X}(\mathcal{F}, \mathcal{F}')$
together with the $G$-linearization
$$
\Theta
:
\sigma^*\mathcal{H}om_{\mathcal{O}_X}(\mathcal{F}, \mathcal{F}')
\xrightarrow{\cong}
\left(\pr_X^{G \times X}\right)^* \mathcal{H}om_{\mathcal{O}_X}(\mathcal{F}, \mathcal{F}')
% ---- 原文行 2060 ----
$$
given by the conjugation
$
\alpha
\mapsto
\theta^{-1}
\comp
\alpha
\comp
\theta'
$
under the canonical isomorphisms\footnote{
These canonical homomorphisms are isomorphisms
because the morphisms
$\sigma$,
$\pr_X^{G\times X}
:
G \times X \rightarrow X$
are flat
(as any group scheme $G$ is flat over the base)
% ---- 原文行 2080 ----
and $\mathcal{F}$ is finitely presented.
}
\begin{align*}
\sigma^* \mathcal{H}om_{\mathcal{O}_X}(\mathcal{F}, \mathcal{F}')
&\xrightarrow{\cong}
\mathcal{H}om_{\mathcal{O}_{G \times X}}\left(\sigma^* \mathcal{F}, \sigma^* \mathcal{F}'\right)
\quad\text{and}\quad
\\[2mm]
\left(\pr_X^{G \times X}\right)^* \mathcal{H}om_{\mathcal{O}_X}(\mathcal{F}, \mathcal{F}')
&\xrightarrow{\cong}
\mathcal{H}om_{\mathcal{O}_{G \times X}}\left(\left(\pr_X^{G \times X}\right)^*\mathcal{F}, \left(\pr_X^{G \times X}\right)^*\mathcal{F}'\right).
\end{align*}

More precisely, for each open $U \subseteq G \times X$ and a local section
$
\alpha \in \left(\sigma^*\mathcal{H}om_{\mathcal{O}_X}(\mathcal{F}, \mathcal{F}')\right)(U)
$
viewed as an $\mathcal{O}_{G\times X}|_U$-module homomorphism
$
\alpha : \sigma^*\mathcal{F}|_U \rightarrow \sigma^*\mathcal{F}'|_U
% ---- 原文行 2100 ----
$,
the section $\Theta(U)(\alpha)$ is defined by the composition
$$
\left(\theta|_U\right)^{-1}
\comp
\alpha
\comp
\theta'|_U
:
\left(\pr_X^{G \times X}\right)^*\mathcal{F}|_U
\rightarrow
\left(\pr_X^{G \times X}\right)^*\mathcal{F}'|_U.
$$

Equipped with the tensor product and internal hom,
$\VB_G(X)$
(as well as $\Coh_G(X)$ if $X$ is locally Noetherian)
forms a \emph{closed symmetric monoidal category}.

\end{defn}
% ---- 原文行 2120 ----
We end this paragraph by computing sections between $G$-equivariant sheaves.
\begin{prop}\label{prop:InvariantSectionsOfEquivariantInternalHom}
Let
$(\mathcal H,\Theta)
:=
\mathcal Hom_{\mathcal O_X}
\bigl((\mathcal F,\theta),(\mathcal F',\theta')\bigr)$.
For every $G$-invariant open subscheme $U\subseteq X$,
we have
$$
\mathcal H(U)^G
=
\Hom_{\Mod_G(U)}
\left(
(\mathcal F,\theta)|_U,\,
(\mathcal F',\theta')|_U
\right).
$$
\end{prop}

% ---- 原文行 2140 ----
\begin{proof}
A section $f\in\mathcal H(U)$ is an $\mathcal O_U$-module
morphism $f:\mathcal F|_U\rightarrow\mathcal F'|_U$.
By the definition of $G$-invariant sections and of $\Theta$,
the following hold:
\begin{align*}
f\in\mathcal H(U)^G
&\quad\Leftrightarrow\quad
\Theta|_{G\times U}(\sigma_U^*f)
=
\left(\pr_U^{G\times U}\right)^*f
\\
&\quad\Leftrightarrow\quad
\theta_U^{-1}\cdot\sigma_U^*f\cdot\theta'_U
=
\left(\pr_U^{G\times U}\right)^*f
\\
&\quad\Leftrightarrow\quad
\sigma_U^*f\cdot\theta'_U
=
% ---- 原文行 2160 ----
\theta_U\cdot
\left(\pr_U^{G\times U}\right)^*f
\end{align*}
The last equality is exactly the $G$-equivariance condition
for $f$.
\end{proof}




